\newpage
\begin{center}

\large\textbf{Юго-Западный государственный университет}
\vskip 1em
\noindent \normalsize {Кафедра \underline{\hspace{3cm}программной инженерии\hspace{5.5cm}}}
%\normalsize{Кафедра программной инженерии}
\vskip 1em
\ifВКР{
        \begin{flushright}
        \begin{tabular}{p{.4\textwidth}}
        \centrow УТВЕРЖДАЮ: \\
        \centrow Заведующий кафедрой \\
        \hrulefill \\
        \setarstrut{\footnotesize}
        \centrow\footnotesize{(подпись, инициалы, фамилия)}\\
        \restorearstrut
        «\underline{\hspace{1cm}}»
        \underline{\hspace{3cm}}
        20\underline{\hspace{1cm}} г.\\
        \end{tabular}
        \end{flushright}
        }\fi
\end{center}
\vspace{1em}
  \begin{center}
  \large
\ifВКР{
ЗАДАНИЕ НА ВЫПУСКНУЮ КВАЛИФИКАЦИОННУЮ РАБОТУ
  ПО ПРОГРАММЕ МАГИСТРАТУРЫ}
  \else
ЗАДАНИЕ НА КУРСОВУЮ РАБОТУ (ПРОЕКТ)
\fi
\normalsize
  \end{center}
\vspace{1em}
{\parindent0pt
  Студента \АвторРод, шифр\ \Шифр, группа \Группа
  
1. Тема «\Тема\ \ТемаВтораяСтрока»
\ifВКР{
утверждена приказом ректора ЮЗГУ от \ДатаПриказа\ № \НомерПриказа
}\fi.

2. Срок предоставления работы к защите \СрокПредоставления

3. Исходные данные для создания программной системы:

3.1. Перечень решаемых задач:}

\renewcommand\labelenumi{\theenumi)}

\begin{enumerate}
\item определить функциональные и технические требования разрабатываемого приложения;
\item разработать концептуальную модель интеллектуальной системы производственного планирования;
\item спроектировать программную систему построения расписаний и адаптивной диспетчеризации;
\item сконструировать и протестировать программную систему с нейросетевыми модулями планирования.
\end{enumerate}


  3.2. Входные данные и требуемые результаты для программы:

\begin{enumerate}
\item Производственные заказы – перечень заказов с указанием сроков сдачи, приоритетов и технологических операций, определяющих последовательность и содержание работ.
\item Производственные ресурсы – сведения о станках и рабочих центрах, их текущем состоянии, коэффициенте готовности, признаках ремонта и специализации по типам операций.
\item Нормативно-справочная информация – перечень типов операций и допустимых связей «ресурс–тип», ограничивающих назначение операций на оборудование.
\item Обучающие выборки – синтетические данные о состояниях оборудования и эвристических назначениях операций, используемые для тренировки нейросетевых моделей.
\item Обученные нейросетевые модели – файлы весов, полученные в результате тренировки моделей-диспетчеров (PointerNet, SLIM, PPO), готовые к загрузке и применению.
\item Производственное расписание – назначенные для каждой операции ресурсы, время старта и окончания, сформированное системой по выбранному методу планирования.
\item Диаграмма Ганта – визуализация построенного расписания с цветовой кодировкой по типам операций и временной шкалой.
\item Рекомендации нейросетевого диспетчера – предлагаемые ресурсы для каждой неразмещённой операции, полученные на основе текущего состояния системы.
\item Метрики обучения – кривые потерь, точность на валидации, матрицы ошибок и отчёты классификации, сохраняемые в процессе тренировки моделей.
\item Сводные отчёты сравнения методов – показатели опоздания, makespan, коэффициент вариации загрузки и другие KPI, вычисляемые при тестировании различных методов планирования.
\end{enumerate}

{\parindent0pt

4. Перечень вопросов, подлежащих исследованию (разработке): }

  4.1  Анализ существующих решений в области производственного планирования и диспетчеризации, выявление их преимуществ и ограничений.

  4.2  Исследование эвристических алгоритмов и методов обучения нейросетевых диспетчеров для задачи гибкого производственного расписания.

  4.3  Сравнительный анализ эффективности нейросетевых и классических методов построения расписаний по ключевым производственным показателям.

{\parindent0pt
5. Перечень графического материала:

\списокПлакатов

\vskip 2em
\begin{tabular}{p{6.8cm}C{3.8cm}C{4.8cm}}
Руководитель \ifВКР{ВКР}\else работы (проекта) \fi & \lhrulefill{\fill} & \fillcenter\Руководитель\\
\setarstrut{\footnotesize}
& \footnotesize{(подпись, дата)} & \footnotesize{(инициалы, фамилия)}\\
\restorearstrut
Задание принял к исполнению & \lhrulefill{\fill} & \fillcenter\Автор\\
\setarstrut{\footnotesize}
& \footnotesize{(подпись, дата)} & \footnotesize{(инициалы, фамилия)}\\
\restorearstrut
\end{tabular}
}

\renewcommand\labelenumi{\theenumi.}
